% GENERATED FILE. Edit canonical lessons, then run npm run generate:books.
% canonical-source-hash: fnv1a64:73d7638226d64813
% canonical-lessons: KA-C399-badatana, KA-C399-karmika, KA-C399-savalattu, KA-C399-vitarisu, KA-C399-gramina

\chapter{Poverty and Need}
\label{ch:ka-poverty-and-need}
\hlchaptermodality{399}

\begin{chapteropening}
\textbf{By the end of this chapter:} \emph{I can say poverty, a labourer, a privilege, to distribute and rural.}

Complete the last lesson of chapter 399: I can say poverty, a labourer, a privilege, to distribute and rural.
\end{chapteropening}

\section[\texorpdfstring{\kn{ಬಡತನ} (baḍatana)}{ಬಡತನ (baḍatana)}]{\kn{ಬಡತನ} (baḍatana) — poverty}

\label{lesson:KA-C399-badatana}



\begin{quote}
Before the new one: say the Kannada for to mobilise, then the Kannada for an activist.
\end{quote}

\subsection*{You'll want to know: \kn{ಬಡತನ}}
\textbf{\kn{ಬಡತನ}} — \emph{baḍatana} — \textquotedblleft{}poverty\textquotedblright{}.

\begin{quote}
\kn{ಬಡತನದಿಂದ} \kn{ಅವನು} \kn{ಶಾಲೆ} \kn{ಬಿಟ್ಟನು}.
\end{quote}

\begin{quote}
\emph{baḍatanadinda avanu śāle biṭṭanu.}
\end{quote}

\begin{quote}
He left school because of poverty.
\end{quote}

\kn{ಬಡ} (poor) + -\kn{ತನ} (-ness).

\begin{cousinweb}
\textbf{\kn{ಬಡತನ}} (\emph{baḍatana}) is native Dravidian: DEDR 5222 lists Kannada \emph{baḍa}, \textquotedblleft{}poor\textquotedblright{}, and \emph{baḍatana}, \textquotedblleft{}poverty\textquotedblright{}, with the ending -\emph{tana}, \textquotedblleft{}-ness\textquotedblright{}.
\end{cousinweb}

\begin{grammarlens}[title={What you've built}]
One more word for giving reasons and opinions.
\end{grammarlens}

\subsection*{Your turn}
\begin{itemize}
\raggedright
  \item \emph{Say it:} \emph{baḍatana}
  \item \emph{Say it:} \emph{baḍatana}, once more
  \item \emph{Say it:} \emph{baḍatanadinda avanu śāle biṭṭanu.} — He left school because of poverty.
  \item \emph{Recall:} say the Kannada for to mobilise, then the Kannada for an activist, then say \emph{baḍatana} again
\end{itemize}

\begin{tcolorbox}[breakable,colback=teal!4,colframe=teal!35!black,title={Before you move on}]
What does \kn{ಬಡತನ} mean? (\textquotedblleft{}Poverty\textquotedblright{}.) Say it once more.
\end{tcolorbox}
\section[\texorpdfstring{\kn{ಕಾರ್ಮಿಕ} (kārmika)}{ಕಾರ್ಮಿಕ (kārmika)}]{\kn{ಕಾರ್ಮಿಕ} (kārmika) — a labourer}

\label{lesson:KA-C399-karmika}



\begin{quote}
Before the new one: say the Kannada for an activist, then the Kannada for poverty.
\end{quote}

\subsection*{You'll want to know: \kn{ಕಾರ್ಮಿಕ}}
\textbf{\kn{ಕಾರ್ಮಿಕ}} — \emph{kārmika} — \textquotedblleft{}a labourer\textquotedblright{}.

\begin{quote}
\kn{ಕಾರ್ಮಿಕರು} \kn{ಹೆಚ್ಚು} \kn{ಸಂಬಳ} \kn{ಕೇಳಿದರು}.
\end{quote}

\begin{quote}
\emph{kārmikaru heccu sambaḷa kēḷidaru.}
\end{quote}

\begin{quote}
The labourers asked for higher pay.
\end{quote}

Plural: \kn{ಕಾರ್ಮಿಕರು}, \emph{kārmikaru}.

\kn{ಕಾರ್ಮಿಕ} \kn{ಸಂಘ}, a trade union.

\begin{cousinweb}
\textbf{\kn{ಕಾರ್ಮಿಕ}} (\emph{kārmika}) is borrowed from Sanskrit \emph{kārmika}, \textquotedblleft{}engaged in action\textquotedblright{} (CDIAL 3076), from \emph{karman}, \textquotedblleft{}work\textquotedblright{}.
\end{cousinweb}

\begin{grammarlens}[title={What you've built}]
One more word for giving reasons and opinions.
\end{grammarlens}

\subsection*{Your turn}
\begin{itemize}
\raggedright
  \item \emph{Say it:} \emph{kārmika}
  \item \emph{Say it:} \emph{kārmika}, once more
  \item \emph{Say it:} \emph{kārmikaru heccu sambaḷa kēḷidaru.} — The labourers asked for higher pay.
  \item \emph{Recall:} say the Kannada for an activist, then the Kannada for poverty, then say \emph{kārmika} again
\end{itemize}

\begin{tcolorbox}[breakable,colback=teal!4,colframe=teal!35!black,title={Before you move on}]
What does \kn{ಕಾರ್ಮಿಕ} mean? (\textquotedblleft{}A labourer\textquotedblright{}.) Say it once more.
\end{tcolorbox}
\section[\texorpdfstring{\kn{ಸವಲತ್ತು} (savalattu)}{ಸವಲತ್ತು (savalattu)}]{\kn{ಸವಲತ್ತು} (savalattu) — a privilege; a benefit}

\label{lesson:KA-C399-savalattu}



\begin{quote}
Before the new one: say the Kannada for poverty, then the Kannada for a labourer.
\end{quote}

\subsection*{You'll want to know: \kn{ಸವಲತ್ತು}}
\textbf{\kn{ಸವಲತ್ತು}} — \emph{savalattu} — \textquotedblleft{}a privilege; a benefit\textquotedblright{}.

\begin{quote}
\kn{ಹಿರಿಯ} \kn{ನಾಗರಿಕರಿಗೆ} \kn{ಬಸ್ಸಿನಲ್ಲಿ} \kn{ಕೆಲವು} \kn{ಸವಲತ್ತುಗಳು} \kn{ಇವೆ}.
\end{quote}

\begin{quote}
\emph{hiriya nāgarikarige bassinalli kelavu savalattugaḷu ive.}
\end{quote}

\begin{quote}
Senior citizens get some concessions on the bus.
\end{quote}

Plural: \kn{ಸವಲತ್ತುಗಳು}, \emph{savalattugaḷu}.

Often plural: \kn{ಸವಲತ್ತುಗಳು}, benefits.

\begin{cousinweb}
\textbf{\kn{ಸವಲತ್ತು}} (\emph{savalattu}) is usually traced to Arabic \emph{suhūlat}, \textquotedblleft{}ease, convenience\textquotedblright{}, which reached Kannada through Hindustani and Marathi.
\end{cousinweb}

\begin{grammarlens}[title={What you've built}]
One more word for giving reasons and opinions.
\end{grammarlens}

\subsection*{Your turn}
\begin{itemize}
\raggedright
  \item \emph{Say it:} \emph{savalattu}
  \item \emph{Say it:} \emph{savalattu}, once more
  \item \emph{Say it:} \emph{hiriya nāgarikarige bassinalli kelavu savalattugaḷu ive.} — Senior citizens get some concessions on the bus.
  \item \emph{Recall:} say the Kannada for poverty, then the Kannada for a labourer, then say \emph{savalattu} again
\end{itemize}

\begin{tcolorbox}[breakable,colback=teal!4,colframe=teal!35!black,title={Before you move on}]
What does \kn{ಸವಲತ್ತು} mean? (\textquotedblleft{}A privilege; a benefit\textquotedblright{}.) Say it once more.
\end{tcolorbox}
\section[\texorpdfstring{\kn{ವಿತರಿಸು} (vitarisu)}{ವಿತರಿಸು (vitarisu)}]{\kn{ವಿತರಿಸು} (vitarisu) — to distribute}

\label{lesson:KA-C399-vitarisu}



\begin{quote}
Before the new one: say the Kannada for a labourer, then the Kannada for a privilege.
\end{quote}

\subsection*{You'll want to know: \kn{ವಿತರಿಸು}}
\textbf{\kn{ವಿತರಿಸು}} — \emph{vitarisu} — \textquotedblleft{}to distribute\textquotedblright{}.

\begin{quote}
\kn{ಶಾಲೆ} \kn{ಮಕ್ಕಳಿಗೆ} \kn{ಉಚಿತ} \kn{ಪುಸ್ತಕಗಳನ್ನು} \kn{ವಿತರಿಸಿತು}.
\end{quote}

\begin{quote}
\emph{śāle makkaḷige ucita pustakagaḷannu vitarisitu.}
\end{quote}

\begin{quote}
The school gave out free books to the children.
\end{quote}

Past (I …): \kn{ವಿತರಿಸಿದೆ}, \emph{vitariside}.

Who receives takes -\kn{ಗೆ}. \kn{ವಿತರಣೆ}, distribution.

\begin{cousinweb}
\textbf{\kn{ವಿತರಿಸು}} (\emph{vitarisu}) is built on Sanskrit \emph{vitaraṇa}, \textquotedblleft{}granting, bestowal, gift\textquotedblright{} (Monier-Williams), with the Kannada verb ending -\emph{isu}.
\end{cousinweb}

\begin{grammarlens}[title={What you've built}]
One more word for giving reasons and opinions.
\end{grammarlens}

\subsection*{Your turn}
\begin{itemize}
\raggedright
  \item \emph{Say it:} \emph{vitarisu}
  \item \emph{Say it:} \emph{vitarisu}, once more
  \item \emph{Say it:} \emph{śāle makkaḷige ucita pustakagaḷannu vitarisitu.} — The school gave out free books to the children.
  \item \emph{Recall:} say the Kannada for a labourer, then the Kannada for a privilege, then say \emph{vitarisu} again
\end{itemize}

\begin{tcolorbox}[breakable,colback=teal!4,colframe=teal!35!black,title={Before you move on}]
What does \kn{ವಿತರಿಸು} mean? (\textquotedblleft{}To distribute\textquotedblright{}.) Say it once more.
\end{tcolorbox}
\section[\texorpdfstring{\kn{ಗ್ರಾಮೀಣ} (grāmīṇa)}{ಗ್ರಾಮೀಣ (grāmīṇa)}]{\kn{ಗ್ರಾಮೀಣ} (grāmīṇa) — rural}

\label{lesson:KA-C399-gramina}



\begin{quote}
Before the new one: say the Kannada for a privilege, then the Kannada for to distribute.
\end{quote}

\subsection*{You'll want to know: \kn{ಗ್ರಾಮೀಣ}}
\textbf{\kn{ಗ್ರಾಮೀಣ}} — \emph{grāmīṇa} — \textquotedblleft{}rural\textquotedblright{}.

\begin{quote}
\kn{ಗ್ರಾಮೀಣ} \kn{ಪ್ರದೇಶಗಳಲ್ಲಿ} \kn{ಆಸ್ಪತ್ರೆಗಳು} \kn{ಕಡಿಮೆ}.
\end{quote}

\begin{quote}
\emph{grāmīṇa pradēśagaḷalli āspatregaḷu kaḍime.}
\end{quote}

\begin{quote}
There are few hospitals in rural areas.
\end{quote}

\kn{ಗ್ರಾಮ}, a village.

\begin{cousinweb}
\textbf{\kn{ಗ್ರಾಮೀಣ}} (\emph{grāmīṇa}) is borrowed from Sanskrit \emph{grāmīṇa}, \textquotedblleft{}of a village, rustic\textquotedblright{} (Monier-Williams), from \emph{grāma}, \textquotedblleft{}a village\textquotedblright{}.
\end{cousinweb}

\begin{grammarlens}[title={What you've built}]
One more word for giving reasons and opinions.
\end{grammarlens}

\subsection*{Your turn}
\begin{itemize}
\raggedright
  \item \emph{Say it:} \emph{grāmīṇa}
  \item \emph{Say it:} \emph{grāmīṇa}, once more
  \item \emph{Say it:} \emph{grāmīṇa pradēśagaḷalli āspatregaḷu kaḍime.} — There are few hospitals in rural areas.
  \item \emph{Recall:} say the Kannada for a privilege, then the Kannada for to distribute, then say \emph{grāmīṇa} again
\end{itemize}

\begin{tcolorbox}[breakable,colback=teal!4,colframe=teal!35!black,title={Before you move on}]
What does \kn{ಗ್ರಾಮೀಣ} mean? (\textquotedblleft{}Rural\textquotedblright{}.) Say it once more.
\end{tcolorbox}
